%% sec/intro.tex — §1 Introduction
%% Source: paper/sec1_intro_v3.md (md remains source of truth)

\section{Introduction}
\label{sec:intro}

Large language models are now active participants in analog sizing
loops, calling SPICE simulators and revising netlists across
iterations~\cite{analogcoder,adollm,eesizer,anaflow,ledro,autosizer,analogsage}.
Reported success rates on open-PDK benchmarks such as
AnalogGym~\cite{analoggym} are encouraging, and
EEsizer~\cite{eesizer} has begun to look at multi-LLM evaluation on
six basic circuits. Two practical questions, however, remain
unanswered as a \emph{conjunction}: \textbf{(i)~which LLM should a
team pick under a token budget when running against a real
industrial PDK?}---no prior work pairs multi-LLM evaluation with
industrial-node deployment, cost--quality Pareto analysis,
\emph{and} per-LLM failure-mode taxonomy; and \textbf{(ii)~is
domestic-China LLM inference industrially viable for analog
sizing?}---to our knowledge zero published work evaluates Moonshot,
MiniMax, Xiaomi, or DeepSeek on this workload.

Concurrent work AnalogAgent~\cite{analogagent} addresses a related
NDA constraint \emph{orthogonally} via on-premise tiny open-weight
models (Qwen3 1.7B--14B); we address the deployment-economics
constraint via an eleven-checkpoint cost--quality Pareto on a real
industrial PDK. The two stories bracket the industrial-deployment
design space.

We present an eleven-LLM benchmark for industrial-node analog
sizing running on \textbf{Safe-Bridge}, a verifiable NDA-safe
wrapper around the PC\,$\leftrightarrow$\,remote-host SKILL IPC
channel that ensures no foundry content reaches any cloud LLM. The
eleven checkpoints span seven vendor families across two countries
(6~US: Anthropic Opus~4.7 / Sonnet~4.6 / Haiku~4.5, OpenAI
GPT-5.5 / GPT-5.4-mini, Google Gemini~2.x; 5~China: Moonshot
Kimi~K2.5, MiniMax~M2.7, Xiaomi MiMo~v2.5-pro, DeepSeek V4-pro /
V4-flash). The workload is LC\_VCO\_20GHz; an OpAmp testbench is
the in-revision hook. The four-vendor China-frontier cohort turns
domestic-China viability from a single anecdote into a cross-vendor
pattern.

\paragraph{Contributions.}
\begin{itemize}
  \item \textbf{SP1.}~Safe-Bridge---the verifiable NDA-safe
    substrate (\S\ref{sec:platform}) that makes SP2 possible on
    real foundry silicon: four code-anchored scrub layers, five
    hard invariants, defense-in-depth audit methodology.
  \item \textbf{SP2 (headline).}~An eleven-LLM industrial-node
    analog-sizing benchmark (\S\ref{sec:methodology}--\S\ref{sec:results}):
    6~US + 5~China across seven vendor families, identical
    Safe-Bridge / spec / prompt; only \texttt{-{}-llm} differs.
  \item \textbf{SP3.}~Cost--quality Pareto
    (\S\ref{sec:results}) with a labelled reasoning-token premium
    sub-section---the sharpest single numerical observation in the
    paper.
  \item \textbf{SP4.}~Failure-mode taxonomy and cross-vendor
    domestic-China LLM viability
    (\S\ref{sec:failures}--\S\ref{sec:discussion}): a
    four-China-frontier-vendor cohort rather than a single point
    estimate.
\end{itemize}

\S\ref{sec:related}--\S\ref{sec:conclusion} cover related work,
Safe-Bridge, methodology, Pareto results, failure taxonomy,
domestic-China discussion, and conclusion respectively.
